\begin{pastexam}[label={px:nchu-115-is-11}]{中興資工所乙組 115 學年度 第 11 題}
（單選；每題 2 分，答錯不倒扣）

Which statement best describes the approach used by dynamic programming?
\begin{enumerate}[label=\Alph*.,nosep]
\item It uses randomization to select among many possible solutions with equal probability.
\item It enumerates all possible solutions and then discards those that violate constraints.
\item It stores solutions to overlapping subproblems and reuses them to avoid redundant
computation.
\item It selects the next step based solely on the immediate local gain without considering
future consequences.
\end{enumerate}
\end{pastexam}

\begin{pastsoln}
答案是 C。

動態規劃適用在有最佳子結構、而且子問題大量重疊的問題（資結書第~18.1~節）。 直接遞迴時同一個
子問題會被算很多次，例如費氏數的遞迴呼叫次數以指數成長（資結書第~18.1.1~節）；把每個子問題
的答案記下來（由上而下的記憶化，或由下而上填表），每個只算一次，時間就變成「子問題個數 $\times$
每個的成本」。
\begin{itemize}
\item A 錯：DP 是確定性的，不靠隨機。
\item B 錯：先列出所有解再刪掉不合格的是窮舉（產生再測試）。 DP 不列出所有解，而是用子問題的
最佳值組出大問題的最佳值。
\item C 對。
\item D 錯：只看眼前的好處、不管後果是貪心法。 DP 正好相反，每一步都比較所有選擇接上子問題
最佳值的結果。
\end{itemize}
\end{pastsoln}
